\documentclass[11pt]{article}
\newcommand{\ListaTitulo}{Lista 13: Simulado Prova 2 (Unidades 3 e 4)}
% Preâmbulo compartilhado: MATF14, Listas de Exercícios 2026
% Usado por ListaNN.tex e GabaritoNN.tex via % Preâmbulo compartilhado: MATF14, Listas de Exercícios 2026
% Usado por ListaNN.tex e GabaritoNN.tex via % Preâmbulo compartilhado: MATF14, Listas de Exercícios 2026
% Usado por ListaNN.tex e GabaritoNN.tex via \input{preamble}

\usepackage[utf8]{inputenc}
\usepackage[T1]{fontenc}
\usepackage[brazil]{babel}
\usepackage{fullpage}
\usepackage{amsmath,amssymb}
\usepackage{enumitem}
\usepackage{xcolor}
\usepackage{fancyhdr}
\usepackage{titlesec}
\usepackage{listings}
\usepackage{hyperref}
\usepackage{booktabs}
\usepackage{graphicx}
\usepackage{tikz}

\definecolor{matf14azul}{HTML}{0D3B54}
\definecolor{matf14dourado}{HTML}{C9971E}

\titleformat{\section}{\normalfont\Large\bfseries\color{matf14azul}}{\thesection}{1em}{}
\titleformat{\subsection}{\normalfont\large\bfseries\color{matf14azul}}{\thesubsection}{1em}{}

\providecommand{\ListaTitulo}{Lista de Exercícios}

\setlength{\headheight}{14pt}
\pagestyle{fancy}
\fancyhf{}
\lhead{\small MATF14: Estatística Econômica I}
\rhead{\small \ListaTitulo}
\cfoot{\thepage}
\renewcommand{\headrulewidth}{0.4pt}

\lstset{
  basicstyle=\ttfamily\small,
  breaklines=true,
  frame=single,
  columns=fullflexible,
  backgroundcolor=\color{gray!8},
  commentstyle=\color{matf14dourado},
  keywordstyle=\color{matf14azul}\bfseries,
  language=R
}

\newcounter{questaocounter}
\newenvironment{questao}[1]{%
  \stepcounter{questaocounter}%
  \bigskip\noindent\textbf{Questão #1.}\ %
}{\par}

\newcommand{\cabecalho}[2]{%
  \begin{center}
    {\Large\bfseries\color{matf14azul} #1}\\[0.3em]
    {\large #2}
  \end{center}
  \vspace{0.5em}
  \noindent\rule{\linewidth}{0.6pt}
  \vspace{0.5em}
}

\newenvironment{solucao}{%
  \par\smallskip\noindent\textit{\color{matf14dourado!80!black}Solução.}\ %
}{\hfill$\blacksquare$\par}


\usepackage[utf8]{inputenc}
\usepackage[T1]{fontenc}
\usepackage[brazil]{babel}
\usepackage{fullpage}
\usepackage{amsmath,amssymb}
\usepackage{enumitem}
\usepackage{xcolor}
\usepackage{fancyhdr}
\usepackage{titlesec}
\usepackage{listings}
\usepackage{hyperref}
\usepackage{booktabs}
\usepackage{graphicx}
\usepackage{tikz}

\definecolor{matf14azul}{HTML}{0D3B54}
\definecolor{matf14dourado}{HTML}{C9971E}

\titleformat{\section}{\normalfont\Large\bfseries\color{matf14azul}}{\thesection}{1em}{}
\titleformat{\subsection}{\normalfont\large\bfseries\color{matf14azul}}{\thesubsection}{1em}{}

\providecommand{\ListaTitulo}{Lista de Exercícios}

\setlength{\headheight}{14pt}
\pagestyle{fancy}
\fancyhf{}
\lhead{\small MATF14: Estatística Econômica I}
\rhead{\small \ListaTitulo}
\cfoot{\thepage}
\renewcommand{\headrulewidth}{0.4pt}

\lstset{
  basicstyle=\ttfamily\small,
  breaklines=true,
  frame=single,
  columns=fullflexible,
  backgroundcolor=\color{gray!8},
  commentstyle=\color{matf14dourado},
  keywordstyle=\color{matf14azul}\bfseries,
  language=R
}

\newcounter{questaocounter}
\newenvironment{questao}[1]{%
  \stepcounter{questaocounter}%
  \bigskip\noindent\textbf{Questão #1.}\ %
}{\par}

\newcommand{\cabecalho}[2]{%
  \begin{center}
    {\Large\bfseries\color{matf14azul} #1}\\[0.3em]
    {\large #2}
  \end{center}
  \vspace{0.5em}
  \noindent\rule{\linewidth}{0.6pt}
  \vspace{0.5em}
}

\newenvironment{solucao}{%
  \par\smallskip\noindent\textit{\color{matf14dourado!80!black}Solução.}\ %
}{\hfill$\blacksquare$\par}


\usepackage[utf8]{inputenc}
\usepackage[T1]{fontenc}
\usepackage[brazil]{babel}
\usepackage{fullpage}
\usepackage{amsmath,amssymb}
\usepackage{enumitem}
\usepackage{xcolor}
\usepackage{fancyhdr}
\usepackage{titlesec}
\usepackage{listings}
\usepackage{hyperref}
\usepackage{booktabs}
\usepackage{graphicx}
\usepackage{tikz}

\definecolor{matf14azul}{HTML}{0D3B54}
\definecolor{matf14dourado}{HTML}{C9971E}

\titleformat{\section}{\normalfont\Large\bfseries\color{matf14azul}}{\thesection}{1em}{}
\titleformat{\subsection}{\normalfont\large\bfseries\color{matf14azul}}{\thesubsection}{1em}{}

\providecommand{\ListaTitulo}{Lista de Exercícios}

\setlength{\headheight}{14pt}
\pagestyle{fancy}
\fancyhf{}
\lhead{\small MATF14: Estatística Econômica I}
\rhead{\small \ListaTitulo}
\cfoot{\thepage}
\renewcommand{\headrulewidth}{0.4pt}

\lstset{
  basicstyle=\ttfamily\small,
  breaklines=true,
  frame=single,
  columns=fullflexible,
  backgroundcolor=\color{gray!8},
  commentstyle=\color{matf14dourado},
  keywordstyle=\color{matf14azul}\bfseries,
  language=R
}

\newcounter{questaocounter}
\newenvironment{questao}[1]{%
  \stepcounter{questaocounter}%
  \bigskip\noindent\textbf{Questão #1.}\ %
}{\par}

\newcommand{\cabecalho}[2]{%
  \begin{center}
    {\Large\bfseries\color{matf14azul} #1}\\[0.3em]
    {\large #2}
  \end{center}
  \vspace{0.5em}
  \noindent\rule{\linewidth}{0.6pt}
  \vspace{0.5em}
}

\newenvironment{solucao}{%
  \par\smallskip\noindent\textit{\color{matf14dourado!80!black}Solução.}\ %
}{\hfill$\blacksquare$\par}


\begin{document}

\cabecalho{Lista de Exercícios 13: Simulado}{Revisão geral, no formato e no nível da Prova 2 (Cap. 3--4, Aulas 29--30)}

\noindent \textit{Simulado de 10 pontos, no mesmo formato da prova oficial. Resolva com tempo
cronometrado (sugestão: 100 minutos, sem consulta) antes de olhar o gabarito em
\texttt{Gabarito13.pdf}.}

\begin{enumerate}

\item Considere uma variável aleatória discreta $X$ com distribuição de probabilidade:
\begin{center}
\begin{tabular}{c|ccccc}
\toprule
$k$ & 1 & 2 & 3 & 4 & 5 \\
\midrule
$P(X=k)$ & $2/10$ & $1/10$ & $3/10$ & $2/10$ & $2/10$ \\
\bottomrule
\end{tabular}
\end{center}
\begin{enumerate}
  \item (1 ponto) Calcule $P(|X-3|\ge 2)$.
  \item (1 ponto) Calcule $E(X)$ e $\mathrm{Var}(X)$.
  \item (1 ponto) Calcule a função de distribuição acumulada de $X$.
\end{enumerate}

\item Uma linha de produção tem 2\% de peças defeituosas, de forma independente entre peças.
Deseja-se a probabilidade de que uma amostra de 150 peças contenha no máximo 3 defeituosas.
\begin{enumerate}
  \item (1 ponto) Calcule essa probabilidade usando a distribuição Binomial exata.
  \item (1 ponto) Calcule a mesma probabilidade usando a aproximação pela distribuição Poisson.
\end{enumerate}

\item Seja $C>0$ uma constante e $X$ uma variável aleatória contínua com densidade
$$
f(x) = \begin{cases} Cx^{-2}, & x\ge 5 \\ 0, & x<5 \end{cases}
$$
\begin{enumerate}
  \item (1 ponto) Encontre o valor de $C$ para que $f$ seja uma densidade de probabilidade válida.
  \item (1 ponto) Calcule a função de distribuição acumulada $F(x)$, para $x\ge5$.
  \item (1 ponto) Calcule $P(X>10)$.
\end{enumerate}

\item O retorno anual de um fundo multimercado é aproximadamente Normal, com média $10\%$ e
desvio padrão $18\%$.
\begin{enumerate}
  \item (1 ponto) Calcule a probabilidade do retorno em um ano ser maior que $15\%$.
  \item (1 ponto) Calcule a probabilidade do retorno em um ano ser menor que $-8\%$.
\end{enumerate}

\end{enumerate}

\vspace{1cm}
\begin{center}
\Large\textbf{Formulário}
\end{center}
\begin{small}
\begin{itemize}
  \item Se $X$ é contínua, $E(X)=\int xf(x)dx$; se discreta, $E(X)=\sum_k k\,P(X=k)$.
    $\mathrm{Var}(X)=E(X^2)-[E(X)]^2$; $dp(X)=\sqrt{\mathrm{Var}(X)}$.
  \item $X\sim\text{Bernoulli}(p)$: $P(X=k)=p^k(1-p)^{1-k}$, $k\in\{0,1\}$; $E(X)=p$,
    $\mathrm{Var}(X)=p(1-p)$.
  \item $X\sim\text{Bin}(n,p)$: $P(X=k)=\binom{n}{k}p^k(1-p)^{n-k}$; $E(X)=np$,
    $\mathrm{Var}(X)=np(1-p)$.
  \item $X\sim\text{Pois}(\lambda)$: $P(X=k)=e^{-\lambda}\lambda^k/k!$; $E(X)=\mathrm{Var}(X)=\lambda$.
  \item $X\sim N(\mu,\sigma^2)$: $E(X)=\mu$, $\mathrm{Var}(X)=\sigma^2$; $Z=(X-\mu)/\sigma\sim
    N(0,1)$.
  \item $X\sim\text{Unif}(\alpha,\beta)$: $f(x)=1/(\beta-\alpha)$; $E(X)=(\alpha+\beta)/2$,
    $\mathrm{Var}(X)=(\beta-\alpha)^2/12$.
  \item $\int x^n\,dx = \dfrac{x^{n+1}}{n+1}+c$ ($n\ne-1$); $\displaystyle\int e^{ax}\,dx =
    \dfrac1a e^{ax}+c$.
\end{itemize}
\end{small}

\end{document}
